% ============================================================
%  Глава 7 — Кватернионы и проекции сенсоров
%  Файл: compute_node/dashboard_node.py
% ============================================================
\chapter{Кватернионы, проекции и преобразования ориентации}
\label{ch:quaternions}

\begin{filebox}[title={Файл исходного кода}]
\filepath{compute\_node/dashboard\_node.py} \quad (Python, строки 298--521)
\end{filebox}

\section{Углы Эйлера из акселерометра (сырые)}

\begin{filebox}[title={\filepath{dashboard\_node.py}, строки 307--308 --- \_mqtt\_imu()}]
\begin{lstlisting}[style=pythonstyle, numbers=left, firstnumber=306]
# Raw pitch/roll from accel (no yaw)
pitch = math.degrees(math.atan2(ay, math.sqrt(ax*ax + az*az)))
roll  = math.degrees(math.atan2(-ax, az))
\end{lstlisting}
\end{filebox}

\textbf{Тангаж} (наклон вперёд/назад):

\begin{equation}
  \theta_{\text{pitch}} = \mathrm{atan2}\!\left(a_y,\, \sqrt{a_x^2 + a_z^2}\right)
  \label{eq:pitch_raw}
\end{equation}

\textbf{Крен} (наклон влево/вправо):

\begin{equation}
  \phi_{\text{roll}} = \mathrm{atan2}(-a_x,\, a_z)
  \label{eq:roll_raw}
\end{equation}

\begin{notebox}
Функция $\mathrm{atan2}(y, x)$ возвращает угол в $(-\pi, \pi]$ и
корректно обрабатывает все четыре квадранта, в отличие от $\arctan(y/x)$.
\end{notebox}

\section{Кватернион $\to$ угол курса (yaw из ROS2 Odometry)}

\begin{filebox}[title={\filepath{dashboard\_node.py}, строки 503--506 --- \_odom\_cb()}]
\begin{lstlisting}[style=pythonstyle, numbers=left, firstnumber=503]
q = msg.pose.pose.orientation
siny_cosp = 2.0 * (q.w * q.z + q.x * q.y)
cosy_cosp = 1.0 - 2.0 * (q.y * q.y + q.z * q.z)
yaw = math.atan2(siny_cosp, cosy_cosp)
\end{lstlisting}
\end{filebox}

\subsection{Кватернионы и ориентация в 3D}

Кватернион единичной длины $\vect{q} = (q_w, q_x, q_y, q_z)$, $|\vect{q}|=1$,
представляет вращение в 3D:

\begin{equation}
  \vect{q} = \cos\frac{\alpha}{2} + \hat{\vect{n}}\sin\frac{\alpha}{2}
  \label{eq:quaternion_def}
\end{equation}

где $\hat{\vect{n}}$ --- ось вращения, $\alpha$ --- угол.

\subsection{Формула извлечения угла курса}

Из матрицы вращения, построенной из кватерниона, угол вращения вокруг $Z$:

\begin{equation}
  R_{32} = 2(q_w q_z + q_x q_y), \quad
  R_{33} = 1 - 2(q_y^2 + q_z^2)
  \label{eq:rot_matrix_elems}
\end{equation}

\begin{equation}
  \boxed{\psi = \mathrm{atan2}\!\bigl(2(q_w q_z + q_x q_y),\;
    1 - 2(q_y^2 + q_z^2)\bigr)}
  \label{eq:yaw_from_quat}
\end{equation}

Это стандартная ZYX-декомпозиция (aerospace convention):
$\psi$ --- курс, $\theta$ --- тангаж, $\phi$ --- крен.

\section{Полярные координаты лазерного сканера $\to$ декартовы}

\begin{filebox}[title={\filepath{dashboard\_node.py}, строки 516--521 --- \_scan\_cb()}]
\begin{lstlisting}[style=pythonstyle, numbers=left, firstnumber=516]
def _scan_cb(self, msg: LaserScan):
    points, angle = [], msg.angle_min
    for r in msg.ranges:
        if msg.range_min < r < msg.range_max:
            points.append([
                round(r * math.cos(angle), 3),
                round(r * math.sin(angle), 3)
            ])
        angle += msg.angle_increment
\end{lstlisting}
\end{filebox}

\textbf{Полярные $\to$ декартовы координаты} (точка облака лазерного сканера):

\begin{equation}
  \begin{pmatrix} x \\ y \end{pmatrix}
  = r \begin{pmatrix} \cos\alpha \\ \sin\alpha \end{pmatrix}
  \label{eq:polar_to_cart}
\end{equation}

где $r$ --- измеренная дальность, $\alpha$ --- угол луча лазера.

\section{Геометрическое представление кватерниона}

\begin{center}
\begin{tikzpicture}[scale=1.2]
  % Coordinate axes
  \draw[-{Stealth}, thick] (0,0,0) -- (2,0,0) node[right]{$X$};
  \draw[-{Stealth}, thick] (0,0,0) -- (0,2,0) node[above]{$Z$};
  \draw[-{Stealth}, thick] (0,0,0) -- (0,0,2) node[right]{$Y$};

  % Rotation axis
  \draw[-{Stealth}, red!70, very thick] (0,0,0) -- (1.2, 1.5, 0.5)
    node[right, red!70]{$\hat{\vect{n}}$};

  % Arc for angle
  \draw[blue!60, thick, ->] (0.6,0,0) arc[start angle=0, end angle=50, radius=0.6]
    node[right, blue!60, font=\small]{$\alpha/2$};

  % Quaternion label
  \node[font=\small] at (0, -0.4, 0) {$\vect{q} = \cos\dfrac{\alpha}{2} + \hat{\vect{n}}\sin\dfrac{\alpha}{2}$};
\end{tikzpicture}
\end{center}

\section{Полная матрица вращения из кватерниона}

\begin{equation}
  \mat{R} = \begin{pmatrix}
    1-2(q_y^2+q_z^2) & 2(q_xq_y-q_wq_z) & 2(q_xq_z+q_wq_y) \\
    2(q_xq_y+q_wq_z) & 1-2(q_x^2+q_z^2) & 2(q_yq_z-q_wq_x) \\
    2(q_xq_z-q_wq_y) & 2(q_yq_z+q_wq_x) & 1-2(q_x^2+q_y^2)
  \end{pmatrix}
  \label{eq:rot_matrix_full}
\end{equation}

Для чистого вращения вокруг $Z$ ($q_x = q_y = 0$,
$q_z = \sin(\psi/2)$, $q_w = \cos(\psi/2)$):

\begin{equation}
  \mat{R}_z(\psi) = \begin{pmatrix}
    \cos\psi & -\sin\psi & 0 \\
    \sin\psi &  \cos\psi & 0 \\
    0 & 0 & 1
  \end{pmatrix}
  \label{eq:Rz_yaw}
\end{equation}
